\paragraph*{Generic \emph{you} and addressed \emph{you}: what the split can and cannot say.}
\vekalet{$M_1$ counts second-person forms and does not ask whom they name. We split the sentences carrying one into \textsc{addressee}, \textsc{generic} and \textsc{undetermined} with a cue list fixed before the measurement, with \textsc{addressee} taking precedence where both fire. Across the 16 models the generic count falls in 16 and the addressed count in 14. A preliminary check by the assistant (an LLM agent, not a human labeller) on 40 sentences found the generic label right $12/13$ of the time and the addressed label $7/14$. In that check the deontic cue (\emph{you should}, \emph{you must}) fired on generic advice as readily as on advice to a reader. Of the $13$ the cues left \textsc{undetermined}, $13$ were resolvable on reading (overall agreement $19/40$). On the authors' labels for $100$ second-person sentences (Appendix~\ref{app:labelcheck}) the classifier labels only $17$. Its generic label is right $5/8$ of the time and its addressed label $6/9$. Of the $83$ it leaves \textsc{undetermined}, the authors read $78$ as addressed or generic. \textbf{The asymmetry of the preliminary check is not reproduced on human labels.} \textsc{Undetermined} is $76\%$ of all second-person sentences in the panel. Both halves are read from the judge-labelled check of Appendix~\ref{app:labelcheck} (\S\ref{sec:measure}). \textbf{We do not claim the withdrawal is confined to either half}: on human labels neither of the classifier's two labels is right often enough to carry that claim.}
